\documentclass[10pt,a4paper]{amsart}
\usepackage[margin=19mm]{geometry}
\usepackage{amsmath,amssymb,mathtools}
\usepackage[scr=rsfs,cal=euler]{mathalfa}
\usepackage{xfrac}
\usepackage[unicode,hidelinks,bookmarks=false]{hyperref}
\newcommand{\ie}{i.e.\ }
\newcommand{\cf}{cf.\ }
\newcommand{\resp}{respectively\ }
\newcommand{\Subsection}{\subsection}
\providecommand{\nobreakdash}{\nobreak\hbox{-}}
\setlength{\emergencystretch}{2em}
\multlinegap0pt
\usepackage{amssymb}
\usepackage{enumerate}
\usepackage{algorithm}\usepackage{algpseudocode}
\usepackage{xcolor}
\usepackage[normalem]{ulem}
\newtheorem{lemma}{Lemma}
\newcommand\Vr{\mathcal{V}}
\newcommand\diam{\mathrm{diam}}
\newcommand\lune{\mathrm{lune}}
\newcommand\rvr{\mathcal{R}}
\title{Computation of degree-1 persistent homology on larger point-clouds using the Reduced Vietoris-Rips filtration}
\author{Musashi Ayrton Koyama, Facundo M\'emoli, Vanessa Robins, Katharine Turner}
\date{Source: arXiv:2307.16333v3 (CC BY 4.0)}
\begin{document}
\maketitle

\noindent\emph{Source and licence.} Computation of degree-1 persistent homology on larger point-clouds using the Reduced Vietoris-Rips filtration. Version arXiv:2307.16333v3 (CC BY 4.0), distributed by arXiv under the Creative Commons Attribution 4.0 International licence, http://creativecommons.org/licenses/by/4.0/, as recorded in the arXiv licence metadata for this exact version and shown on the abstract page https://arxiv.org/abs/2307.16333v3. Full licence notice: \texttt{licenses/arxiv-2307.16333-CC-BY-4.0.txt}.\par\smallskip

\noindent\textit{Editorial scope.} Counted unit: the article's lemma q-the-big-lemma (source0.tex lines 2286--2289). The article's complete original proof of it is delivered with it (source0.tex lines 2291--2366, one \texttt{proof} environment of 76 lines). No mathematics is written, repaired or re-derived: the statement and the proof are the article's own lines, in the article's own notation, and no retained line is substituted or re-flowed.  Retained uncounted as the source-local context the proof needs: lines 2242--2245; lines 2258--2261; lines 2269--2281.  Nothing was omitted from any retained span, so the package carries no omission record.  No retained line contains a citation, so the wrapper carries no bibliography and no bibliography entry.  No retained line contains a figure environment, an image-inclusion command or drawing code, so no image is delivered and the package declares no asset dependency.  Editorial formatting only, in addition: an AMS \texttt{amsart} preamble that restates the article's own declarations and macros used by the retained lines, namely the environments \texttt{lemma} on separate counters, together with the standard packages those lines need.  The retained environments print in this document as Lemma 1 in order of appearance, whereas the article prints the same items with the numbering of its own full document.  The complete native source of the article as distributed in the arXiv e-print for this exact version is bundled unchanged as \texttt{sources/144976/source0.tex}.
\begin{lemma}
\label{q-theta-is-surjective}
    $\theta_{r}$ is surjective for all $r \geq 0$. 
\end{lemma}

\begin{lemma}
\label{q-isomorphism-inclusions}
    Let $s > 0$. If $\theta_{s}$ is an isomorphism, then it follows that $B_{q}(\Vr_{s}(X)) = B_{q}(\rvr^{q}_{s}(X))$. 
\end{lemma}

\begin{lemma}
\label{q-first-edge-proof}
    Consider $r>0$ arbitrary and suppose that $\theta_{s}$ is injective for all $s<r$. Let all the $q$-simplices of diameter $r$ be $\tau_{1}  < \cdots < \tau_{m_r} $. Then $\lune (\tau_{1}) = \emptyset$ and thus
    \begin{equation}
    \label{q-same-edges-equation-1}
         \{ \partial (\langle  x \tau_{1} \rangle) \; | \; x\in \lune (\tau_{1}) \} = \emptyset 
    \end{equation}
    and thus we trivially have 
    \begin{equation}
    \label{q-same-edges-equation-1-1}
         \{ \partial (\langle  x \tau_{1} \rangle) \; | \; x\in \lune (\tau_{1}) \} \subset B_{q}(\rvr_{r}^{q}(X)).
    \end{equation}
\end{lemma}

\begin{lemma}
\label{q-the-big-lemma}
    Consider $r>0$ arbitrary and suppose that $\theta_{s}$ is injective for all $s<r$. Let $\sigma$ be a $(q+1)$-simplex in $\Vr_{r}(X)$ such that $\diam (\sigma) = r$, then we have that $\partial \sigma \in B_{q}(\rvr_{r}^{q}(X))$. That is we have $B_{q}(\rvr_{r}^{q}(X)) = B_{q}(\Vr_{r}(X))$.
\end{lemma}

\begin{proof}
    First observe that by Lemma \ref{q-theta-is-surjective} that $\theta_s$ is an isomorphism for all $s < r$. This means that $B_{q}(\rvr_{s}^{q}(X)) = B_{q}(\Vr_{s}(X))$ by Lemma \ref{q-isomorphism-inclusions}. 
    
    Let all the $q$-simplices of diameter $r$ be $ \tau_{1}  < \cdots <  \tau_{m_{r}} $ and denote $\tau_{k} = \langle y_{0k}....,y_{qk} \rangle$. We wish to show that all $(q+1)$-simplices $\sigma$, with $\diam(\sigma) = r$ are such that $\partial \sigma \in B_{q}(\rvr_{r}^{q}(X))$. This is equivalent to showing that 

    \begin{equation}
    \label{q-same-edges-equation}
         \{ \partial ( \langle  x \tau_{j} \rangle)\; | \; x\in \lune (\tau_{j}) \} \subset B_{q}(\rvr_{r}^{q}(X))
    \end{equation}

for $j = 1,...,m_{r}$. We will prove that (\ref{q-same-edges-equation}) holds for $j = 1,...,m_{r}$ by induction. That is we will show the following:

\begin{itemize}
    \item That (\ref{q-same-edges-equation}) holds for $j=1$. This is just Lemma \ref{q-first-edge-proof}. 

    \item If (\ref{q-same-edges-equation}) holds for $j=1,...,k$, then it holds for $j = k+1$. 
\end{itemize}



We now show that if (\ref{q-same-edges-equation}) holds for $j = 1,...,k$, then it holds for $j = k+1$. 

Let $x$ be any point in $\lune (\langle y_{0 (k+1)}...y_{q (k+1)} \rangle)$ and consider the \\ $(q+1)$-simplex $\langle x y_{0(k+1)}...y_{q(k+1)} \rangle$. Let $w$ be the point in the connected component of $\lune (\langle y_{0(k+1)}...y_{q(k+1)} \rangle)$ that also contains $x$ such that $w\in L(\langle y_{0(k+1)}...y_{q(k+1)}\rangle)$. If $x=w$ then there is nothing to prove since $\sigma \in \rvr^{q}_{r}(X)$ by definition and thus $\partial \sigma \in B_{1}(\rvr_{r}^{q}(X))$. Otherwise, since $x$ and $w$ are in the same component we know there exists $v_{0},...,v_{t} \in \lune (\langle y_{0(k+1)}...y_{q(k+1)} \rangle)$ such that $x = v_0, v_1, ..., v_t = w$ forms a path with

\begin{equation}
\langle v_{i}v_{i+1}y_{0(k+1)}...\hat{y}_{l_1(k+1)}...\hat{y}_{l_2(k+1)}...y_{q(k+1)} \rangle < \langle y_{0(k+1)}...y_{q(k+1)} \rangle 
\label{path-equation}
\end{equation}

for all $0 \leq l_1 < l_2 \leq q$.

We know that $\partial (\partial (\langle y_{0(k+1)}...y_{q(k+1)} v_{i}v_{i+1} \rangle )) = 0$ and thus 
\begin{equation}
\begin{aligned}
    \partial (\langle y_{0(k+1)}...y_{q(k+1)} v_{i} \rangle ) = 
    \partial (\langle y_{0(k+1)}...y_{q(k+1)} v_{i+1} \rangle) \\ + \sum_{l=0}^{q}\partial (\langle v_{i}v_{i+1}y_{0(k+1)}...\hat{y}_{l(k+1)}...y_{q(k+1)} \rangle)
    \end{aligned}
    \label{q-equation-general-proof}
\end{equation}
We now show that $\partial (\langle v_{i}v_{i+1}y_{0(k+1)}...\hat{y}_{l(k+1)}...y_{q(k+1)} \rangle) \in B_{q}(\rvr_{r}^{q}(X))$ for $l = 0,...,q$ :

\begin{itemize}
    \item $\langle v_{i}v_{i+1}y_{0(k+1)}...\hat{y}_{l_{1}(k+1)}...\hat{y}_{l_{2}(k+2)}...y_{q(k+1)} \rangle < \langle y_{0(k+1)}...y_{q(k+1)} \rangle$ \\ for all $0 \leq l_1 < l_2 \leq q$ by (\ref{path-equation}). 

    \item $\langle v_{i} y_{0(k+1)}...\hat{y}_{l(k+1)}...y_{q(k+1)} \rangle < \langle y_{0(k+1)}...y_{q(k+1)} \rangle $ for $l =0,...,q$ since 
    
    \vspace{3pt}
    $v_{i} \in \lune (\langle y_{0(k+1)}...y_{q(k+1)} \rangle)$. 

    \item $\langle v_{i+1} y_{0(k+1)}...\hat{y}_{l(k+1)}...y_{q(k+1)} \rangle < \langle y_{0(k+1)}...y_{q(k+1)} \rangle $ for $l =0,...,q$ since 
    
    \vspace{3pt}
    $v_{i+1} \in \lune (\langle y_{0(k+1)}...y_{q(k+1)} \rangle)$.
    
\end{itemize}

For $l= 0,...,q$, all faces of $\langle v_{i}v_{i+1}y_{0(k+1)}...\hat{y}_{l(k+1)}...y_{q(k+1)} \rangle$ appear in the filtration before $\langle y_{0(k+1)}...y_{q(k+1)} \rangle$. For each $l=0,...,q$ we either have all \\faces of $\langle v_{i}v_{i+1}y_{0(k+1)}...\hat{y}_{l(k+1)}...y_{q(k+1)} \rangle$ have diameter less than $r$ or at least one of them has diameter $r$. If they all have diameter less than $r$ then \\ $\partial (\langle v_{i}v_{i+1}y_{0(k+1)}...\hat{y}_{l(k+1)}...y_{q(k+1)} \rangle) \in B_{q}(\Vr_{s}(X)) = B_{q}(\rvr^{q}_{s}(X))$ for $s < r$. If at \\ least one of the faces of $\langle v_{i}v_{i+1}y_{0(k+1)}...\hat{y}_{l(k+1)}...y_{q(k+1)} \rangle$ has diameter $r$ then let $\eta$ be the face with diameter $r$ that appears latest in the filtration. Let $z$ be the \\ point in $\langle v_{i}v_{i+1}y_{0(k+1)}...\hat{y}_{l(k+1)}...y_{q(k+1)} \rangle$ that does not appear in $\eta$, then we have that $z\in \lune(\eta)$. Thus $\partial (\langle v_{i}v_{i+1}y_{0(k+1)}...\hat{y}_{l(k+1)}...y_{q(k+1)} \rangle) = \partial(\langle z \eta\rangle)$. Since $\eta$ \\ comes before $\tau_{k+1}$ in the filtration it follows that 

\begin{equation}
\partial (\langle v_{i}v_{i+1}y_{0(k+1)}...\hat{y}_{l(k+1)}...y_{q(k+1)} \rangle) \in \{\partial (\langle x \tau_{j} \rangle) | x \in \lune (\tau_{j})\}
\end{equation}

for some $j = 1,...,k$.
Plugging $i = t-1$ into  (\ref{q-equation-general-proof}) we obtain the following.

\begin{equation}
\begin{aligned}
\label{q-equation-general-proof-q-1}
    \partial (\langle y_{0(k+1)}...y_{q(k+1)} v_{t-1} \rangle ) = 
    \partial (\langle y_{0(k+1)}...y_{q(k+1)} v_{t} \rangle)\\ + \sum_{l=0}^{q}\partial (\langle v_{t-1}v_{t}y_{0(k+1)}...\hat{y}_{l(k+1)}...y_{q(k+1)} \rangle)
    \end{aligned}
\end{equation}
We have already shown that the summation on the right hand side of (\ref{q-equation-general-proof-q-1}) is in $B_{q}(\rvr_{r}^{q}(X))$. Remembering that $v_{t} = w$, the fact that the first term on the right hand side of (\ref{q-equation-general-proof-q-1}) is in $B_{q}(\rvr_{r}^{q}(X))$ follows from $\langle v_{t}\tau_{k+1}\rangle$ being a $(q+1)$-simplex in $\rvr_{r}^{q}(X)$ by the definition of $\rvr_{r}^{q}(X)$. Thus it follows that $\partial (\langle v_{t-1} \tau_{k+1} \rangle) \in B_{q}(\rvr_{r}^{q}(X))$. Using (\ref{q-equation-general-proof}) repeatedly by letting $i = t-2, t-3, ...,0$ we can subsequently show that $\partial (\langle \tau_{k+1}v_{t-2}),...,\partial(\langle \tau_{k+1}v_{0}\rangle) \in B_{q}(\rvr_{r}^{q}(X))$. Remembering that $v_{0} = x$ we thus have that $\partial (\langle x\tau_{k+1} \rangle)   \in B_{q}(\rvr_{r}^{q}(X))$. 


\end{proof}

\end{document}
